\section{Threat Model and the Estimand}\label{sec:estimand}

\subsection{Threat model}\label{sec:threat}

Two roles recur. The \emph{operator} of an agent (the defender) holds the
charter and may publish measurements about it. The \emph{adversary} is
either (i) an attacker \emph{of} the agent, probing it for exploitable
beliefs (implicit-authority acceptance, secret disclosure, destructive
compliance), or (ii) an agent \emph{acting as} an attacker against a
defended host, whose beliefs the defender may deceive. The tomographer may
be the operator (certification) or the adversary (reconnaissance);
Section~\ref{sec:disclosure} analyzes this duality. All interactions are black-box: only channel behavior (transcripts and
tool-call traces) is observed. The adversary controls one thing only: the
content of an interaction episode---the message, the tool output, the timing.
It controls no weights, no system prompt, and no provider filter, and observes
nothing but the response, censored under the same rules as any other caller
(Section~\ref{sec:e2}). Three things are out of scope by construction: the
agent's internal representations, the training-time provenance of a displayed
belief, and the intent behind its charter. One capability is \emph{in} scope
deliberately: the target may be evaluation-aware; Corollary~\ref{cor:pub}
exists because a finite public suite cannot rule that out. Transfer assumes
the target cannot condition on probe \emph{provenance} beyond probe content
(assumption (A2)); targets that can void transfer by construction.

\paragraph{Attack classes and probes}
The probe families instantiate, in instrument form, the attack classes
documented for production LLM agents---indirect prompt injection
\cite{greshake2023,liu2024promptinjection}, its tool-integrated and dynamic
benchmarks \cite{injecagent2024,agentdojo2024}. Those works ask whether an
attack \emph{succeeds};
tomography asks what the response reveals about the agent's beliefs. Every
experiment runs on locally constructed foils, on author-administered
infrastructure, or---for the field layers---on commercial endpoints reached
through authenticated, metered API access (Section~\ref{sec:limits}).

\begin{definition}[Agent, charter, benign operation]\label{def:agent}
An \emph{agent} is a (possibly randomized) map from interaction histories to
responses, $A:\mathcal{H}\to\Delta(\mathcal{R})$, with response kernel
$\kappa_B(h,\cdot)$ parameterized by a belief hypothesis $B\in\mathfrak{B}$. Its
\emph{charter} $\chi$ is a published set of policy claims. \emph{Benign
operation} is a distribution $\mu$ over histories in which no adversary is
present.
\end{definition}

\begin{definition}[Belief proposition and probe family]\label{def:probe}
A \emph{belief proposition} $\pi$ is a testable statement about the agent's
operative beliefs (e.g., ``believes forged system authority''). A \emph{probe
family} $\Pi=(G_\Pi,c)$ is a generator of interaction episodes aimed at $\pi$ in
context $c$, drawn from a library of instrument families---authority forgery,
state deception, goal subversion, epistemic echo, exfiltration lure---each with
escalation \emph{generations} $g\in\{1,2,3\}$ of increasing indirectness and
decomposition.
\end{definition}

\begin{definition}[Behavioral belief strength]\label{def:beta}
Given a binary behavioral-consistency functional $X_B(\Pi)\in\{0,1\}$ (the
verdict that the agent's response is consistent with holding $\pi$, produced by
deterministic oracles and a calibrated judge), the \emph{behavioral belief
strength} is
\[
\beta(\pi;\Pi)=\E_{G_\Pi,\kappa_B}\!\left[X_B(\Pi)\right]\in[0,1].
\]
The \emph{belief profile} is $\boldsymbol{\beta}=\{\beta(\pi;\Pi_c)\}_{\pi\in\mathcal{P}}$
over a proposition catalogue $\mathcal{P}$ and context grid $c$---the
tomogram; a bare $B$ indexes a belief hypothesis of
Section~\ref{sec:t1} except in $\Delta(B)$ (Definition~\ref{def:gain}) and
$B_{\mathrm{pre}},B_{\mathrm{post}}$ (Section~\ref{sec:e3}), where it
indexes the measured attacker or profile.
\end{definition}

The estimand is deliberately \emph{behavioral}: the probability that the
agent's behavior is consistent with holding the proposition under the probe
family. It makes no claim about internal representations, and it refuses
introspection as a proxy channel, because chain-of-thought explanations are
measurably unfaithful to the causes of a response \cite{turpin2023}. Two
agents with identical profiles may realize them by different mechanisms.
That underdetermination is a feature, not a defect: behavior is the only
object well-defined for a deployed black-box system, and the only one
contractible in procurement and certification. Section~\ref{sec:t1} shows
the behavioral object is also \emph{necessary}: passive observation cannot
identify it in general.

\paragraph{Relation to attack-success-rate reporting}
On one proposition under one probe family, $\beta$ reduces to the
attack-success rate benchmark suites report
\cite{agentdojo2024,injecagent2024}. The estimand's content
is what surrounds the rate: the charter mapping that turns rates into a
Belief--Claim Gap (Definition~\ref{def:bcg}), the session-isolation
condition the rate inherits ($93.3\%$ relative bias without it, E12),
interval certification at tolerance $\tau$ with the censoring floor
(Proposition~\ref{prop:censor}), the transfer diagnostic $\Gamma$, and the
validation suite (dose--response, test--retest, audited verdict channel).

\begin{definition}[Belief--Claim Gap]\label{def:bcg}
For a charter claim $\varphi\in\chi$ with instrument mapping
$F(\varphi)\subseteq\mathcal{P}\times\mathcal{C}$ (the propositions and contexts
that test it), the \emph{Belief--Claim Gap} is
\[
\mathrm{BCG}(\varphi)=\frac{1}{|F(\varphi)|}\sum_{(\pi,c)\in F(\varphi)}
\beta(\pi;\Pi_c).
\]
$\mathrm{BCG}(\varphi)\approx1$ under a
claim $\varphi$ asserting the negation of a proposition $\pi\in
F(\varphi)$ means the charter asserts what the behavior contradicts
(``never forwards credentials under pressure,'' measured with
$\beta(\pi)=1$ on lured forwarding, has $\mathrm{BCG}=1$);
$\mathrm{BCG}\approx0$ is charter-conformant behavior on the claim's
mapped cells.
\end{definition}

\begin{definition}[Epistemic distance and deception gain]\label{def:gain}
For two profiles $\boldsymbol{\beta},\boldsymbol{\beta}'$, let $d$ be
the per-proposition shift collection
$\{|\beta(\pi)-\beta'(\pi)|\}_{\pi}$ aggregated by the deception
operation's objective weighting---the convention the field estimator of
Section~\ref{sec:field} reports (per-proposition shifts first, then the
objective aggregation). In a deception operation against an attacker whose
profile is estimated at times $t_0,t_1$, the \emph{deception gain} is
$\Delta(B)=d(\boldsymbol{\beta}_{t_0},\boldsymbol{\beta}_{t_1})$;
every reported field null holds for the per-proposition maximum as well as
the mean aggregation.
\end{definition}
